\section{Noisy and Lossy Simulations}
\label{sec:noisy}

\subsection{Loss Model}

Photon loss is modeled as a fictitious beamsplitter of transmissivity $\eta \in [0,1]$ on each mode before detection. The lossy interferometer is $U_{\mathrm{loss}} = \sqrt{\eta} U \oplus \sqrt{1-\eta} I_{\mathrm{env}}$, where the environment modes are traced out. For threshold detection, loss reduces the effective photon number and introduces vacuum contributions.

\subsection{Distinguishability}

Partial distinguishability is modeled by a distinguishability parameter $\xi \in [0,1]$, where $\xi=1$ is fully indistinguishable and $\xi=0$ is fully distinguishable. The output probability is a convex combination:
\begin{equation}
p_{\xi}(\mathbf{n}) = \xi \, p_{\mathrm{bos}}(\mathbf{n}) + (1-\xi) \, p_{\mathrm{dist}}(\mathbf{n}).
\end{equation}
This interpolates between bosonic interference and classical multinomial routing.

\subsection{Preliminary Noise Results}

We evaluated the optimized $n=8$ path circuit under loss ($\eta=0.9$) and partial distinguishability ($\xi=0.8$):

\begin{table}[ht]
\centering
\caption{Noise robustness of optimized $n=8$ path circuit.}
\label{tab:noise_robustness}
\begin{tabular}{lcccc}
\toprule
Condition & $P_{\mathrm{CF}}$ & $P_{\mathrm{success}}$ & $P_{\mathrm{opt|CF}}$ & $N_{\mathrm{shots}}$ \\
\midrule
Ideal & 0.733 & 0.537 & 0.732 & 1.86 \\
Loss $\eta=0.9$ & 0.582 & 0.381 & 0.655 & 2.62 \\
Distinguishability $\xi=0.8$ & 0.691 & 0.482 & 0.697 & 2.07 \\
Both $\eta=0.9, \xi=0.8$ & 0.543 & 0.332 & 0.611 & 3.01 \\
\bottomrule
\end{tabular}
\end{table}

Loss reduces $P_{\mathrm{CF}}$ more than $P_{\mathrm{opt|CF}}$, as expected since loss introduces vacuum and lower-photon-number terms that are not collision-free. Partial distinguishability reduces the interference contrast, lowering both $P_{\mathrm{CF}}$ and $P_{\mathrm{opt|CF}}$.

\subsection{Noise-Aware Objective}

For hardware with known loss $\eta$ and distinguishability $\xi$, one can optimize the noise-aware objective:
\begin{equation}
P_{\mathrm{success}}(\eta,\xi) = \sum_{\mathbf{n}\in\mathrm{CF}, C(\mathbf{n})=C^*} p_{\eta,\xi}(\mathbf{n}).
\end{equation}
Preliminary results show that optimizing directly for the noisy objective yields circuits more robust to the specific noise model, at the cost of reduced ideal performance.

\subsection{Limitations of Current Noise Model}

Our noise model includes only uniform loss and uniform distinguishability. Real devices exhibit:
\begin{itemize}
\item Mode-dependent loss and crosstalk
\item Phase noise and drift
\item Dark counts and detector inefficiency
\item Higher-order photon-number correlations
\item Non-Gaussian noise
\end{itemize}
A full device-specific noise model requires calibration data from the target QPU. The QPU protocol (Section~\ref{sec:qpu}) includes loss and distinguishability characterization as mandatory calibration steps.
